% =============================================================================
% Template: theorem-uc.tex
% Skill:    iacr-math-prose
%
% Use for: a theorem + proof in the UC simulation-based style (Canetti 2001,
%          revised 2020 ePrint 2000/067). Defines ideal functionality F,
%          real protocol pi, simulator S; shows IDEAL ~_c REAL for every
%          environment Z.
%
% Invariants (do not delete):
%   - F is in a labelled figure environment (framed box), not in prose.
%   - pi is in a construction block, with one phase per party role.
%   - S is in pseudocode or a construction block.
%   - Corruption model is explicit (static / adaptive, semi-honest / malicious).
%   - Hybrid model is explicit (e.g., F_CRS-hybrid, F_RO-hybrid).
%   - UC framework citation: Canetti 2001 (FOCS) or Canetti 2020 (ePrint 2000/067).
%   - The indistinguishability argument reduces to a stated assumption.
%
% See:
%   - references/proof-style-uc.md (worked example: commitment from CRS)
%   - references/theorem-statement-style.md
% =============================================================================

% --- 1. Ideal functionality ---------------------------------------------------
\begin{figure}[t]
\begin{framed}
\textbf{Functionality $\TODO{\F_\textsf{Target}}$.}
$\TODO{\F_\textsf{Target}}$ proceeds as follows, parameterised by parties
$\TODO{P_1, \ldots, P_n}$ and adversary $\Sim$.
\begin{itemize}
  \item \textbf{Inputs.}
        \TODO{State precisely what messages F accepts and from whom.}
  \item \textbf{Internal state.}
        \TODO{State the recorded data, session map, corruption set.}
  \item \textbf{Outputs.}
        \TODO{State precisely what messages F emits and to whom.}
  \item \textbf{Corruption.}
        \TODO{State the corruption model and simulator interface.}
\end{itemize}
\end{framed}
\caption{\TODO{Ideal functionality F.}}
\label{fig:\TODO{f-target}}
\end{figure}

% --- 2. Real protocol ---------------------------------------------------------
\begin{construction}[\TODO{$\pi_\textsf{Target}$ in the hybrid model}]
\label{constr:\TODO{pi-target}}
Let \TODO{state the cryptographic primitives invoked.}

\textbf{Setup.}
\TODO{State what the hybrid functionalities provide.}

\textbf{Phase 1 (party $P_i$).}
\begin{itemize}
  \item \TODO{Step 1.}
  \item \TODO{Step 2.}
\end{itemize}

\textbf{Phase 2 (party $P_j$).}
\begin{itemize}
  \item \TODO{Step 1.}
\end{itemize}
\end{construction}

% --- 3. Simulator -------------------------------------------------------------
\paragraph{Simulator $\Sim$.}
On behalf of the corrupted parties, $\Sim$ proceeds as follows.
\begin{itemize}
  \item \TODO{Setup-phase action.}
  \item \textbf{On message from $\F_\textsf{Target}$.} \TODO{Action.}
  \item \textbf{On message from corrupted party (i.e., from $Z$).} \TODO{Action.}
\end{itemize}

% --- 4. Theorem and proof -----------------------------------------------------
\begin{theorem}[UC security of $\TODO{\pi_\textsf{Target}}$]
\label{thm:\TODO{pi-target-uc}}
In the $\TODO{\F_\textsf{Hybrid}}$-hybrid model, the protocol
$\TODO{\pi_\textsf{Target}}$ of \cref{constr:\TODO{pi-target}} UC-realises the
ideal functionality $\TODO{\F_\textsf{Target}}$ (\cref{fig:\TODO{f-target}})
against a \TODO{static, malicious} PPT adversary corrupting any
\TODO{strict subset of the parties}, assuming \TODO{stated assumption}.
\end{theorem}

\begin{proof}[Proof sketch]
Fix a PPT environment $Z$ and a PPT adversary $A$. We show
$\textsf{IDEAL}_{\TODO{\F_\textsf{Target}}, \Sim, Z}
\approxc
\textsf{REAL}_{\TODO{\pi_\textsf{Target}}, A, Z}$
in the $\TODO{\F_\textsf{Hybrid}}$-hybrid model.

\paragraph{Case 1: \TODO{first corruption pattern}.}
\TODO{Describe $Z$'s REAL-world view.}
\TODO{Describe $Z$'s IDEAL-world view as simulated by $\Sim$.}
If $Z$ distinguishes the two views with non-negligible probability $\delta$,
we build a PPT \TODO{assumption-breaker} $B$ with
\[
  \TODO{\Adv}^{\TODO{\textsf{ASSUMPTION}}}_{B}(\secparam) \;\geq\; \delta,
\]
contradicting the assumption.

\paragraph{Case 2: \TODO{second corruption pattern}.}
\TODO{Argument; often a perfect or statistical indistinguishability case.}

Combining the cases,
\[
  \left| \Pr[Z(\textsf{REAL}) = 1] - \Pr[Z(\textsf{IDEAL}) = 1] \right|
  \;\leq\; \TODO{\Adv^{\textsf{ASSUMPTION}}_{B}(\secparam) + \negl(\secparam)}.
  \qed
\]
\end{proof}

% Cite the UC framework version explicitly in the surrounding text:
%
%   "We work in the UC framework of Canetti
%   (FOCS 2001, latest revision ePrint 2000/067)."
%
% For GUC (global UC) extensions, cite Canetti--Dodis--Pass--Walfish, TCC 2007.
